\subsection{The integral source through every modulus prime to six}

\begin{statement}{E.52}{Proved}{The complete congruence source and its lower-central residual}
Let $B=C_2*C_3=\langle s,r\mid s^2=r^3=1\rangle$,
let $F=\ker(B\to\mathfrak S_3)$, and use the actual
generators $x=(sr)^{-2}$ and $y=sxs$.
Let $N=\ker(B\to\operatorname{SL}_2(\mathbf Z/12)/\{\pm I\})$
with the matrix map of E.50.4, and put $H=F'F^3$.
We write hats for their profinite completions, identified
with the corresponding open subgroups of $\widehat B$.
All commutator and power subgroups in a profinite group
are closed. Set
\[
 A=\mathbf Z[1/6],\qquad
 R=\prod_{p\geq5}\mathbf Z_p,\qquad
 K_2=\overline{\widehat F'\widehat F^2},\qquad
 R_{\mathrm{nil}}=\bigcap_{j\geq2}\gamma_j\widehat H.
 \tag{E.52.1}
\]
For every integer $m>1$ prime to six, the integral
eight-pair construction below defines a specified
continuous epimorphism
\[
 \rho_m:\widehat B\twoheadrightarrow G_m
   \subset\operatorname{Sp}_4(\mathbf Z/m)^{48}
                         \rtimes\mathfrak S_3.
 \tag{E.52.2}
\]
The alternating form and the right matrix convention
are specified below. Reduction for $m\mid n$ gives
the transition map $G_n\twoheadrightarrow G_m$,
with $\rho_m$ equal to reduction of $\rho_n$.
Let $G_R$ be their inverse limit and let
$\rho_R:\widehat B\twoheadrightarrow G_R$ be the
resulting continuous map.

There are explicit injective homomorphisms
\[
 \begin{split}
 \Phi_m:\operatorname{SL}_2(\mathbf Z/m)
                    \times(\mathbf Z/m)^5&\hookrightarrow G_m,\\
 \Phi_R:\operatorname{SL}_2(R)\times R^5&\hookrightarrow G_R,
 \end{split}
 \tag{E.52.3}
\]
compatible with every reduction. For $T=\mathbf Z/m$
or $T=R$, write $\rho_T,\Phi_T,G_T$ for these maps.
The rank-three direct summand $V_3(T)\subset T^5$
is the row span of the following displayed matrix:
\[
 D=\begin{pmatrix}
 -4/3&4/3&-2/3&5/3&1\\
 1/3&-1/3&-1/3&-2/3&0\\
 2/3&-2/3&1/3&-4/3&-1
 \end{pmatrix}.
 \tag{E.52.4}
\]
Then the exact images, with the inclusions of
\textup{(E.52.3)}, are
\[
 \begin{aligned}
 \rho_T(\widehat N)=\rho_T(K_2)
     &=\Phi_T\bigl(\operatorname{SL}_2(T)\times T^5\bigr),\\
 \rho_T(\widehat H)&=\rho_T(\widehat F),\\
 C_T:=\rho_T(\widehat F')=\rho_T(\widehat H')
     &=\Phi_T\bigl(\operatorname{SL}_2(T)\times V_3(T)\bigr),\\
 \rho_T(\gamma_j\widehat H)=\rho_T(R_{\mathrm{nil}})
     &=C_T\qquad(j\geq2),\\
 \rho_T(\widehat N')=C_T'
     &=\Phi_T\bigl(\operatorname{SL}_2(T)\times0\bigr).
 \end{aligned}
 \tag{E.52.5}
\]
In particular every element of $C_R$ has a preimage in
the single closed subgroup $R_{\mathrm{nil}}$.
For one such preimage $w$, its entire fibre there is
$w(R_{\mathrm{nil}}\cap\ker\rho_R)$.

The actual finite quotient maps from $\widehat B$ give
\[
 \begin{gathered}
 G_m/\rho_m(K_2)\simeq A_2,\qquad |A_2|=24,\\
 G_m/C_m\simeq A_{2m},\qquad
 A_n=\widehat B/\overline{\widehat F'\widehat F^n},\\
 |G_m|=24m^5|\operatorname{SL}_2(\mathbf Z/m)|,\qquad
 \rho_m(\widehat F)^{\mathrm{ab}}\simeq(\mathbf Z/2m)^2.
 \end{gathered}
 \tag{E.52.6}
\]
The last isomorphism sends $x,y$ to the two ordinary
exponent-sum basis vectors. The corresponding quotients
over $R$ are the inverse limits of these same maps.
Reduction at $m=5,7$ recovers the specified sources
of E.50 and E.51 through their retained orbit
intertwiners. Thus \textup{(E.52.5)} includes all their
prime-power and simultaneous congruence refinements.
\end{statement}

\begin{proof}
We first construct the integral maps and then prove
surjectivity and every asserted quotient using those maps.

\paragraph{Integral matrices and the actual induced source.}
Let $J$ be the four-by-four alternating matrix with
$J_{i,i+1}=1$, $J_{i+1,i}=-1$ and all other entries
zero, and put $\Omega=J^{-1}$.
The matrices used with ordinary multiplication are
\[
 M_i=(I+e_i e_i^{\mathsf t}J)^{\mathsf t}
       \quad(1\leq i\leq4),\qquad
 M_i^{\mathsf t}\Omega M_i=\Omega.
 \tag{E.52.7}
\]
This is the transpose convention which realizes the
right-permutation convention of E.50 and E.51.
Their entries are integers. Direct multiplication
gives the adjacent Artin relations, commutation at
distance at least two, and
\[
 (M_1M_2M_3M_4)^5
       =(M_1^2M_2M_3M_4)^4=-I.
 \tag{E.52.8}
\]
Define $A_{ij}$ from these $M_i$ by the literal pure
braid word
\[
 A_{ij}=M_{j-1}\cdots M_{i+1}M_i^2
                          M_{i+1}^{-1}\cdots M_{j-1}^{-1}.
 \tag{E.52.9}
\]
The first five evaluation pairs are, in order,
\[
 \begin{gathered}
 (A_{12},A_{23}A_{24}),\quad
 (A_{13}A_{23},A_{34}),\quad(A_{23},A_{34}),\\
 (A_{12}A_{13},A_{24}A_{34}),\quad(A_{12},A_{23}).
 \end{gathered}
 \tag{E.52.10}
\]
The final three are
$(M_{i-1}\cdots M_1M_1\cdots M_{i-1},M_i^2)$,
for $i=2,3,4$.
Each pair defines a homomorphism from the discrete
free group $F=\langle x,y\rangle$ to
$\operatorname{Sp}_4(A,\Omega)$.

Order the six matrices of $\operatorname{SL}_2(\mathbf F_2)$
lexicographically. For each, take the first positive
$s,r$ word given by breadth-first search starting at
the identity, trying $s$ before $r$.
These are the representatives $t_q$ in the certificate.
For $b\in B$, the component at $(q,j)$ of the induced
map is the $j$-th pair evaluated at
\[
                       t_q b t_{q\bar b}^{-1}\in F.
 \tag{E.52.11}
\]
The top component is right multiplication by $\bar b$
on the six cosets. The equality
$t_qbb't_{q\bar b\bar b'}^{-1}
 =(t_qbt_{q\bar b}^{-1})(t_{q\bar b}b't_{q\bar b\bar b'}^{-1})$
proves the homomorphism identity, with this exact action
on the 48 coordinates. Reducing the matrices modulo
$m$ and extending from the finitely generated discrete
source to its profinite completion gives
\textup{(E.52.2)}. Here $G_m$ denotes its actual image.
The image in the inverse limit is closed and surjects
at each finite stage, so it is the whole inverse limit.

On the identity-coset third pair, rows and columns two
and three give a literal two-by-two quotient:
the corresponding rows have zero entries in columns
one and four. Therefore extraction defines a
homomorphism on the image of $F$, with
\[
 \theta(x)=\begin{pmatrix}1&0\\2&1\end{pmatrix},\qquad
 \theta(y)=\begin{pmatrix}1&-2\\0&1\end{pmatrix}.
 \tag{E.52.12}
\]
The same extraction defines its reductions and the
continuous map over $R$.

\paragraph{The finite exact factorization over $A$.}
We describe precisely the rational matrix certificate
used here. The rank-97 basis $n_0,\ldots,n_{96}$ of
$N$ is obtained from the spanning tree of the finite
$(2,3)$ coset graph of
$\operatorname{SL}_2(\mathbf Z/12)/\{\pm I\}$.
The program \texttt{de\_ferrari\_a6\_source\_data.py}
retains all edges and the actual free words and
compares these words with the earlier saved source.
It checks both inverse substitutions between the
free graph basis and $x,y$ before any evaluation.
For each of the 48 coordinates, solve over $\mathbf Q$
the finite intertwiner equations
\[
             M_j(n_i)T_j=T_j\theta(n_i)
                     \quad(0\leq i\leq96).
 \tag{E.52.13}
\]
There are 24 coordinates with one-dimensional
intertwiner space and 24 with zero space.
In each nonzero case choose the primitive integral
matrix given by the saved nullspace basis and set
\[
 P_j=(T_j^{\mathsf t}\Omega T_j)^{-1}T_j^{\mathsf t}\Omega,
 \quad P_jT_j=I_2,
 \quad L_j(B)=I_4+T_j(B-I_2)P_j.
 \tag{E.52.14}
\]
For a zero space put $L_j(B)=I_4$.
All entries of $P_j$ lie in $A$; their denominators
divide six. Multiplication and $P_jT_j=I_2$ prove
$L_j(B)L_j(C)=L_j(BC)$.
The formula is the specified action on the symplectic
plane $\operatorname{im}T_j$ and the identity on its
orthogonal complement. Thus it preserves $\Omega$
for every $B\in\operatorname{SL}_2$ over every $A$-algebra.

Flatten the 48 matrices
$L_j(\theta(n_i))^{-1}M_j(n_i)-I_4$ into a row
with 768 entries. The first five rows are a basis
$\mathcal D_0,\ldots,\mathcal D_4$; write
$D_{u,j}$ for the four-by-four piece of $\mathcal D_u$
at coordinate $j$. The exact checks give
\[
 \begin{gathered}
 D_{u,j}D_{v,j}=0,
 \qquad D_{u,j}^{\mathsf t}\Omega+\Omega D_{u,j}=0,\\
 D_{u,j}L_j(B)=L_j(B)D_{u,j},\qquad
 M_j(n_i)=L_j(\theta(n_i))
                    \left(I_4+\sum_{u=0}^4a_{iu}D_{u,j}\right).
 \end{gathered}
 \tag{E.52.15}
\]
For commutation it suffices to check the two elementary
one-parameter matrix coefficients; the independent audit
checks the elementary matrices and their linear
coefficients. The identities involving the form can
also be recovered from the symplectic matrices
$I+D_{u,j}$ and $I-D_{u,j}$: their difference gives
the linear identity, since two is invertible, and
$D_{u,j}D_{v,j}=0$ then kills every quadratic term.
All these are exact rational identities, with no
finite-field rank substituted for a rational rank.

The basis minor in flattened columns
$4,6,12,20,48$ (numbered from zero) has determinant
$73728=2^{13}3^2$.
The matrix $(6a_{iu})$ is integral and its row lattice
has Smith invariants $1,6,6,6,6$.
It follows that the $\mathcal D_u$ remain linearly
independent over every $A$-algebra and the 97
coefficient rows span $A^5$.
The certificate retains all 4656 original integer
matrices, the 48 intertwiner systems, all $T_j,P_j$,
all five global basis rows, and all 97 coefficients.
Each equality in \textup{(E.52.15)} is verified by
literal matrix multiplication.

Define the homomorphism of \textup{(E.52.3)} by
\[
 \Phi_T(B,a)_j
       =L_j(B)\left(I_4+\sum_{u=0}^4a_uD_{u,j}\right),
                  \quad\hbox{with top component }1.
 \tag{E.52.16}
\]
The identities \textup{(E.52.15)} prove multiplicativity
and preservation of $\Omega$.
At coordinate 18, the identity-coset third pair, the
local radical is zero and extraction in
\textup{(E.52.12)} gives exactly $B$.
After $B$ is known, the displayed unit minor recovers
$a$. This proves injectivity and explicitly describes
the inverse on its image.

\paragraph{Surjectivity on $N$ and the whole even-exponent subgroup.}
The elements $x^6,y^6$ belong to $N$ and have
$\theta$-images $E_{21}(12)$ and $E_{12}(-12)$.
For finite $T=\mathbf Z/m$, twelve is a unit and
their powers give every upper and lower elementary
matrix. These generate $\operatorname{SL}_2(T)$:
over each local factor a unimodular column has a
unit entry, elementary elimination reduces it to
$(u,0)$, and the diagonal matrix is a product of
elementary matrices. The Chinese remainder maps
give the assertion over $T$.
This group is perfect. Indeed with
$h=\operatorname{diag}(2,1/2)$,
\[
 [h,E_{12}(t)]=E_{12}(3t),\qquad
 [h,E_{21}(t)]=E_{21}(-3t/4).
 \tag{E.52.17}
\]
Both coefficients are units, so every elementary
generator belongs to the commutator subgroup.

The coefficient rows above also prove that the
additive projection of $\rho_m(N)$ is all of $T^5$.
Its image under \textup{(E.52.16)} is therefore a
subdirect product of the perfect group
$\operatorname{SL}_2(T)$ and the abelian group $T^5$.
Goursat's lemma makes its common quotient both
perfect and abelian, hence trivial: the commutator
subgroup of an abelian group is the identity.
Thus it is the whole displayed direct product.
In particular \textup{(E.52.16)} has image in $G_m$.

The five actual Schreier generators of $K_2$ in $F$
are
\[
 x^2,\quad yxy^{-1}x^{-1},\quad y^2,
       \quad xyxy^{-1},\quad xy^2x^{-1}.
 \tag{E.52.18}
\]
Their 240 induced matrices all have the same
factorization \textup{(E.52.16)}. In the five radical
coordinates, their rows are respectively
\[
 \begin{pmatrix}
 -1/2&1/3&-2/3&2/3&1/2\\
 4/3&-4/3&2/3&-5/3&-1\\
 1/3&0&0&0&0\\
 -5/6&2/3&-1/3&4/3&1/2\\
 -1/3&2/3&-1/3&1/3&0
 \end{pmatrix}.
 \tag{E.52.19}
\]
Also $N\subset K_2$: reduction of the actual modular
map modulo four reads the two exponent parities on
$F$, and $N$ is killed by that map.
Hence $\rho_m(N)=\rho_m(K_2)$ is the whole group in
the first line of \textup{(E.52.5)}.
Every statement here is compatible with reduction.
Compactness of $\widehat N$ and $K_2$ now gives the
same equality over $R$, by taking the closed fibres
over all finite moduli. It also proves topological
perfectness of $\operatorname{SL}_2(R)$ from the
perfectness of every finite image.

\paragraph{Detection of the complete quotient $A_2$.}
For each of the 24 nonzero intertwiner coordinates,
direct multiplication gives signs
$\varepsilon_{x,j},\varepsilon_{y,j}\in\{1,-1\}$
such that
\[
 M_j(x)T_j=\varepsilon_{x,j}T_j\theta(x),\qquad
 M_j(y)T_j=\varepsilon_{y,j}T_j\theta(y).
 \tag{E.52.20}
\]
At coordinate 10 these signs are $(1,-1)$; at
coordinate 26 they are $(-1,1)$.
They are literal distinct coordinates of the
six-coset construction, and both matrices have
$P_jT_j=I_2$ over $A$.
For any $f\in\widehat F$, multiplication extends
\textup{(E.52.20)} to the two corresponding parity
characters of its exponent sums.
On an element of the image of $K_2$, formula
\textup{(E.52.16)} has sign $+1$ at every coordinate;
the radicals annihilate these standard planes.
As two is invertible and $\theta(f)$ is invertible,
\textup{(E.52.20)} forces both parities to vanish
whenever $\rho_m(f)\in\rho_m(K_2)$.
Thus the full inverse image of $\rho_m(K_2)$ in
$\widehat B$ is exactly $K_2$: the top permutation
first forces membership in $\widehat F$, and these
two signs force membership in $K_2$.
This proves the first quotient of
\textup{(E.52.6)}, and its order is
$[B:F][F:F'F^2]=6\cdot4=24$.

\paragraph{The commutator image and every lower-central term.}
Conjugating every $n_i$ by $s,r,x,y$ gives another
literal factorization in \textup{(E.52.16)}.
The new coefficient rows equal the old row multiplied
by a specified five-by-five matrix. All 388 rows
are checked. Since the coefficient projection on
$N$ is a homomorphism, equality on its 97 free
generators proves the equality for every element
of $N$, at every finite modulus and over $R$.
The matrices for $x$ and $y$ commute and square to
the identity. This also follows directly from
$\rho_T(K_2)$ being the displayed direct product,
where conjugation acts trivially on its additive
factor.

The closed normal subgroup $\widehat F'$ is normally
generated by $c=[x,y]$. Its additive image is therefore
generated by the four conjugates with conjugators
$1,x,y,xy$, since $K_2$ acts trivially on that
additive factor. The first three rows are exactly
the rows of \textup{(E.52.4)}; the fourth is
$(1/3,-1/3,2/3,1/3,0)$.
The first three are independent, their minor in columns
$0,2,3$ (numbered from zero) has determinant $-1/3$, and
the four rows have coefficient-lattice Smith
invariants $1,1,1$ in that basis.
They generate precisely $V_3(T)$ at every modulus.
Moreover $\theta(\widehat F')$ is the full
$\operatorname{SL}_2(T)$, because $\theta(\widehat F)$
is onto and this target is perfect.
The same perfect--abelian Goursat argument proves
the formula for $C_T$.

In the ordered basis \textup{(E.52.4)}, conjugation
by $x$ and $y$ acts on row vectors on the right by
\[
 X=\begin{pmatrix}0&1&0\\1&0&0\\-1&-1&-1\end{pmatrix},
 \qquad
 Y=\begin{pmatrix}0&0&1\\-1&-1&-1\\1&0&0\end{pmatrix}.
 \tag{E.52.21}
\]
The row lattice of $X-I$ and $Y-I$ has Smith
invariants $1,1,4$. Since four is a unit in $A$,
\[
 [\rho_T(\widehat F),C_T]=C_T.
 \tag{E.52.22}
\]
Indeed these differences give the whole additive
factor, while commutators inside the perfect
$\operatorname{SL}_2(T)$ factor give that factor.
The equality $C_T'=\Phi_T(\operatorname{SL}_2(T)\times0)$,
and the same equality for $\rho_T(\widehat N')$,
follow directly from the two direct products.

For finite $T$, the index of $C_T$ in the image of
$F$ is $4m^2$, by the already proved group orders.
The elements $x^2,y^2$ in \textup{(E.52.19)} have
$\theta$-images $E_{21}(4),E_{12}(-4)$; these have
order dividing $m$, and their additive coordinates
are also killed by $m$.
Therefore $x^{2m}=y^{2m}=1$ in $G_m$.
The abelianization of $\rho_m(\widehat F)$ is generated
by these two elements, each of order dividing $2m$,
and has order $4m^2$. The actual exponent-sum map
therefore identifies it with $(\mathbf Z/2m)^2$.
Its full kernel in $\widehat F$ is
$\overline{\widehat F'\widehat F^{2m}}$.
Retaining the same $s,r$ and the quotient to
$\mathfrak S_3$ proves $G_m/C_m\simeq A_{2m}$ as
the actual quotient of $\widehat B$.
This establishes all orders and maps in
\textup{(E.52.6)}.

Multiplication by three is invertible on
$(\mathbf Z/2m)^2$. Since $H=F'F^3$, its image
is consequently the whole image of $F$.
Thus its derived image is $C_T$, and
\textup{(E.52.22)} proves inductively the image
$C_T$ for every $\gamma_j\widehat H$, $j\geq2$.
The statements over $R$ follow through the same
compatible finite maps and compactness.

Finally fix $c\in C_R$. The closed sets
\[
 \gamma_j\widehat H\ \cap\
      \rho_m^{-1}(c\bmod m),\qquad j\geq2,\quad(m,6)=1,
 \tag{E.52.23}
\]
are nonempty by the proved finite images.
Any finite collection contains the set obtained
by taking the largest $j$ and the least common
multiple of its moduli, so has nonempty intersection.
Compactness of $\widehat H$ gives a common element
of all these sets. It belongs to $R_{\mathrm{nil}}$
and maps to the prescribed $c$.
The inverse inclusion follows from its containment
in $\widehat H'$. This proves the residual equality
and, by the elementary fibre formula for a group
homomorphism, the stated full word coset.
\end{proof}

The two root programs
\texttt{e\_symplectic\_five\_integral\_probe\_root.py}
and
\texttt{e\_symplectic\_five\_integral\_residual\_root.py}
retain the exact matrices, source words, coordinate
maps and integer lattices used in this proof.
Their recorded runs take 4.643 and 11.403 seconds.
The first program also compares all sixteen pair
permutations literally with the frozen models at
five and seven. The independent integral audit
reconstructs all 4656 $N$ matrices and all 240
$K_2$ matrices and verifies the elementary standard
maps and the radical products.
The second independent audit replays the complete residual
certificate and checks all five $K_2$ Schreier generators
under each of $s,r,x,y$, retaining the row-action convention.
A further literal comparison reconstructs the induced $s,r$
permutations at five and seven: all 29952 and 115200
entries, respectively, agree with the frozen uncompressed
sources. The coset sections agree, so the intervening
conjugating gauge is the identity.
This result supplies the complete congruence source
and its residual image. The common-source and full
peripheral-fibre comparisons with the other retained
quotients are separate maps to be established below;
the four global tasks remain open at this stage.
